%% ====================================================================
\section{Marco teórico}
\label{sec:marco}
%% ====================================================================

Este capítulo desarrolla el soporte teórico que en el informe intermedio quedaba
solamente esbozado. La idea central es leer la coordinación distribuida como un
problema de \emph{clearing} de capacidad efectiva. En el caso homogéneo, la unidad que
se intercambia es un robot; en el caso heterogéneo, la unidad relevante es
\[
  s_{ik}=c_i\,\Psi(T_k(p_i)),
\]
es decir, la capacidad nominal del robot $i$ descontada por la accesibilidad real a la
tarea $k$. Esta lectura conecta el modelo con juegos poblacionales, juegos potenciales,
consenso dinámico y control multi-robot \citep{sandholm2010population,
quijano2017population,barreiro2017distributed,kia2019tutorial,bullo2009distributed}.

El capítulo usa una convención de estados para no mezclar resultados de distinta
madurez:
\begin{itemize}
  \item \textbf{[PROBADO+GATE]}: resultado demostrado y verificado por un script o gate
  numérico del repositorio.
  \item \textbf{[PROBADO]}: resultado demostrado formalmente, aunque sin gate directo o
  con verificación indirecta.
  \item \textbf{[CONDICIONADO]}: resultado correcto bajo una hipótesis técnica que debe
  cerrarse antes de elevarlo a teorema principal.
  \item \textbf{[FUTURO]}: extensión de frontera, útil para la memoria final o una línea
  doctoral, pero no para afirmar como garantía cerrada del TFM.
\end{itemize}

Dicho de forma compacta, la narrativa que debe sostener el TFM es
\[
  \text{cota de capacidad}
  \rightarrow y_i\in\simplex
  \rightarrow s_{ik}
  \rightarrow \Zeff_k
  \rightarrow F_{ik}
  \rightarrow V
  \rightarrow \text{Smith}
  \rightarrow E
  \rightarrow \text{validación}.
\]
La parte probada vive en el núcleo homogéneo y en el modelo fluido agregado; la
implementación con DAC, uniciclo, obstáculos, modos y saturaciones se interpreta como
una perturbación que debe cerrarse con resultados experimentales.

\subsection{Mapa conceptual: consenso, juego, precio y campo}

La arquitectura del método puede resumirse en cuatro verbos: estimar, valorar, decidir
y moverse. El consenso estima ocupación o capacidad; el juego valora cada tarea; Smith
decide cómo mover masa decisional; y el campo convierte la decisión en movimiento de
AGV. La Figura~\ref{fig:marco_capas} muestra el bucle conceptual.

\begin{figure}[H]
\centering
\begin{tikzpicture}[
  block/.style={draw=VIUDark,fill=VIULight,rounded corners=2pt,
    align=center,minimum width=3.1cm,minimum height=0.9cm,font=\small},
  arrow/.style={-{Latex[length=2mm]},thick,color=VIUOrange!90!black}
]
\node[block] (world) {mundo local\\pose, vecinos, tareas};
\node[block, right=0.85cm of world] (cons) {consenso\\$\hat{x}_{ik}$, $\hat{C}_{ik}$};
\node[block, right=0.85cm of cons] (pay) {payoff\\$f_{ik}$, $F_{ik}$};
\node[block, below=0.85cm of pay] (smith) {Smith\\$p_i$ o $y_i$};
\node[block, below=0.85cm of cons] (field) {campo\\$\nabla E$, obstáculos};
\node[block, below=0.85cm of world] (robot) {robot físico\\$v_i,\omega_i$};
\draw[arrow] (world) -- (cons);
\draw[arrow] (cons) -- (pay);
\draw[arrow] (pay) -- (smith);
\draw[arrow] (smith) -- (field);
\draw[arrow] (field) -- (robot);
\draw[arrow] (robot) -- (world);
\draw[arrow] (smith.east) to[out=0,in=0,looseness=1.2] (pay.east);
\end{tikzpicture}
\caption{Lectura del lazo cerrado. El consenso estima hechos; el juego decide papeles;
el campo mueve ruedas; y el movimiento cambia otra vez los hechos disponibles.}
\label{fig:marco_capas}
\end{figure}

Este bucle no es una arquitectura jerárquica de asignador central, planificador global y
controlador local. La propuesta se acerca más a la tradición de control distribuido en
redes robóticas \citep{bullo2009distributed,ren2008distributed}: cada agente ejecuta la
misma ley con información parcial y el comportamiento útil debe aparecer como propiedad
colectiva.

\subsection{Juegos poblacionales y simplex de decisión}

En un juego poblacional, una masa de agentes se reparte entre estrategias. En este TFM,
la estrategia $k$ representa comprometerse con la tarea $k$, y la estrategia $0$
representa inactividad. Para cada robot $i$ se mantiene una preferencia mixta
\[
  y_i=(y_{i0},y_{i1},\ldots,y_{iK})\in\simplex^{K+1},
  \qquad
  \simplex^{K+1}=\left\{y\ge0:\sum_{k=0}^K y_k=1\right\}.
\]
La estrategia pura ejecutada por el robot puede obtenerse como
$s_i=\argmax_k y_{ik}$, con histéresis para evitar chattering. La variable mixta permite
razonar con ecuaciones diferenciales; la variable pura permite ejecutar coaliciones
enteras.

\begin{definicion}[Ocupación, ocupación efectiva y capacidad efectiva]
En el caso homogéneo se define la ocupación fluida de la tarea $k$ como
\[
  z_k=\sum_{i=1}^N y_{ik}.
\]
Con descuento espacial se usa la ocupación efectiva
\[
  \zeff_k=\sum_{i=1}^N \Psi(T_k(p_i))\,y_{ik}.
\]
En el caso heterogéneo la magnitud de clearing es la capacidad efectiva
\[
  \Zeff_k=\sum_{i=1}^N c_i\Psi(T_k(p_i))\,y_{ik}.
\]
\end{definicion}

La diferencia conceptual es importante. En el informe intermedio bastaba hablar de
cardinalidad mínima $|\mathcal{C}_k|\ge n_k$. En el reporte avanzado, la cardinalidad
aparece como caso particular de una condición física más general:
\[
  \sum_{i\in\mathcal{C}_k} c_i\ge M_k.
\]

\subsection{Payoff como presión de reclutamiento}

El payoff homogéneo del informe intermedio tiene la forma
\[
  f_{ik}=r_k\,\Phi(N\hat{x}_{ik},n_k)\,\Psi(d_{ik}),
  \qquad
  f_{i0}=r_0.
\]
La sigmoide
\[
  \Phi(z,n)=\frac{1}{1+\exp(\beta(z-n))}
\]
es decreciente: si faltan robots, la tarea paga más; si sobran, paga menos. Esto invierte
la lectura clásica de muchos juegos de congestión. Aquí no se penaliza únicamente la
saturación; se recluta para cubrir una cota inferior de servicio.

\begin{observacion}[Cota inferior frente a capacidad máxima]
En muchos modelos de asignación o congestión, el problema es evitar que demasiados
agentes usen el mismo recurso. En transporte cooperativo de cargas pesadas ocurre lo
contrario: una tarea falla si recibe menos robots de los necesarios. Esa inversión es
la fuente del payoff sigmoidal y de la regla de \emph{water-filling}.
\end{observacion}

En capacidad heterogénea el payoff estructuralmente integrable es
\begin{equation}
  F_{ik}=c_i\,r_k\,\Phi_k(\Zeff_k)\,\Psi(T_k(p_i)),
  \qquad F_{i0}=r_0.
  \label{eq:marco_payoff_hetero}
\end{equation}
La capacidad aparece dos veces: dentro de $\Zeff_k$ porque contribuye al estado de la
tarea, y fuera porque escala la productividad marginal del robot. Si se elimina uno de
los dos lugares, se pierde la simetría de derivadas cruzadas que permite construir un
potencial.

\subsection{Dinámicas de Smith}

La dinámica de Smith transfiere masa desde estrategias con menor payoff hacia
estrategias con mayor payoff:
\begin{equation}
  \dot y_{ik}
  =
  \sum_{\ell=0}^{K} y_{i\ell}\pos{F_{ik}-F_{i\ell}}
  -
  y_{ik}\sum_{\ell=0}^{K}\pos{F_{i\ell}-F_{ik}}.
  \label{eq:marco_smith}
\end{equation}
Su ventaja operativa frente al replicador es que compara estrategias por pares, sin
necesitar una media global exacta del payoff \citep{sandholm2010population,
quijano2017population}. Eso encaja con grafos móviles: el robot ya calcula el vector de
payoffs local, pero no tiene por qué conocer una media poblacional fiable.

\begin{lema}[Invariancia del simplex]
Si $y_i(0)\in\simplex^{K+1}$, entonces la solución de
\eqref{eq:marco_smith} satisface $y_i(t)\in\simplex^{K+1}$ para todo $t\ge0$.
\end{lema}

\begin{proof}
Si $y_{ik}=0$, el término de salida de la estrategia $k$ se anula y el término de
entrada es no negativo; por tanto ninguna componente cruza a valores negativos. Además,
al sumar \eqref{eq:marco_smith} en $k$ se obtiene
\[
  \sum_k\dot y_{ik}
  =
  \sum_{k,\ell}y_{i\ell}\pos{F_{ik}-F_{i\ell}}
  -
  \sum_{k,\ell}y_{ik}\pos{F_{i\ell}-F_{ik}}=0,
\]
porque los dos sumatorios son iguales tras intercambiar $k$ y $\ell$. La masa total se
conserva.
\end{proof}

\subsection{Juegos potenciales y teorema homogéneo}

El caso homogéneo con $\Psi\equiv1$ admite un potencial exacto. Sea
\[
  z_k=\sum_i y_{ik},
  \qquad
  G_k(z)=\int_0^z\Phi_k(s)\dd s.
\]
El potencial de utilidad es
\begin{equation}
  V(y)=r_0\sum_i y_{i0}+\sum_{k=1}^{K}r_kG_k(z_k).
  \label{eq:marco_potencial_hom}
\end{equation}
Como
\[
  \frac{\partial V}{\partial y_{ik}}
  =
  r_kG_k'(z_k)
  =
  r_k\Phi_k(z_k),
\]
el juego es de potencial exacto en el sentido de \citet{monderer1996potential}.

\begin{proposicion}[Smith asciende el potencial homogéneo] \textbf{[PROBADO]}
Para el potencial \eqref{eq:marco_potencial_hom}, la dinámica de Smith satisface
$\dot V\ge0$. La igualdad sólo puede mantenerse cuando no existe masa colocada en una
estrategia estrictamente peor que otra disponible.
\end{proposicion}

\begin{proof}
Usando $\partial V/\partial y_{ik}=F_{ik}$,
\[
  \dot V=\sum_{i,k}F_{ik}\dot y_{ik}.
\]
Al sustituir \eqref{eq:marco_smith} y agrupar por pares $(k,\ell)$ se obtiene una suma
de términos no negativos proporcional a
\[
  \sum_{i,k,\ell}y_{i\ell}\pos{F_{ik}-F_{i\ell}}^2.
\]
Por tanto $V$ no decrece. Si existe masa en una estrategia $\ell$ y otra estrategia
$k$ paga estrictamente más, el término correspondiente es positivo; en equilibrio no
puede persistir esa desigualdad dentro del soporte.
\end{proof}

\begin{observacion}[Escala de curvatura]
La curvatura de la sigmoide debe leerse con la variable correcta. En variables de
ocupación $z_k$,
\[
  \frac{\dd^2}{\dd z_k^2}\,r_kG_k(z_k)
  =
  r_k\Phi_k'(z_k),
  \qquad
  r_k\Phi_k'(n_k)=-\frac{r_k\beta}{4}.
\]
Si se trabaja con fracciones $x_k=z_k/N$, aparece el factor de escala de la
normalización escogida para el potencial. Esos factores no cambian la definitud
negativa de la Hessiana restringida, pero sí la cota de paso discreto.
\end{observacion}

\subsection{Regla de water-filling}

El equilibrio fluido se caracteriza por una igualación de payoff marginal. Para una
tarea activa,
\[
  r_k\Phi_k(z_k^*)=\lambda^*.
\]
La inversa sólo debe evaluarse dentro de su dominio. Para evitar que la notación
$[\cdot]_+$ oculte logaritmos fuera de rango, se define la demanda elemental como
\begin{equation}
  W_k(\lambda)=
  \begin{cases}
  0,
  & \lambda\ge r_k\Phi_k(0),\\[4pt]
  n_k+\dfrac{1}{\beta}\log\!\left(\dfrac{r_k}{\lambda}-1\right),
  & 0<\lambda<r_k\Phi_k(0).
  \end{cases}
  \label{eq:marco_water_hom}
\end{equation}
Para una tarea activa, $z_k^*=W_k(\lambda^*)$.

\begin{teorema}[Water-filling homogéneo] \textbf{[PROBADO+GATE]}
En el caso homogéneo, el equilibrio fluido satisface \eqref{eq:marco_water_hom}. Si hay
abundancia de robots, $\lambda^*=r_0$. Si hay escasez, $\lambda^*$ se determina por
\[
  \sum_{k=1}^{K}W_k(\lambda^*)=N.
\]
\end{teorema}

\begin{proof}
En un equilibrio de Nash interior, dos tareas activas no pueden ofrecer payoff marginal
distinto, porque Smith movería masa desde la tarea peor pagada hacia la mejor pagada.
Por tanto todas las tareas activas igualan un nivel $\lambda^*$. Si queda masa inactiva,
el pago de inactividad fija ese nivel en $r_0$. Si no queda masa inactiva, la restricción
de masa total fija $\lambda^*$ mediante $\sum_k z_k^*=N$.
\end{proof}

\begin{corolario}[Puente entero]
Si
\[
  r_k\ge (1+e^\beta)\lambda^*,
\]
entonces el equilibrio fluido cumple $z_k^*\ge n_k+1$. Este margen explica por qué el
clearing entero no es un detalle de implementación, sino el puente entre el modelo
continuo y coaliciones físicas.
\end{corolario}

\begin{proof}
En la rama activa de \eqref{eq:marco_water_hom}, la condición
$z_k^*\ge n_k+1$ equivale a
\[
  \frac{1}{\beta}\log\left(\frac{r_k}{\lambda^*}-1\right)\ge1.
\]
Al exponentiar y despejar se obtiene $r_k\ge(1+e^\beta)\lambda^*$.
\end{proof}

\subsection{Subasta evolutiva de un reloj}

El precio de tarea puede adaptarse con una ley de error de servicio:
\begin{equation}
  \dot r_k=\gamma\,[q_k-\Zeff_k(r)],
  \qquad r_k\in[r_0,r_{\max}].
  \label{eq:marco_subasta}
\end{equation}
Esta ley no añade un subastador central. Cada tarea ajusta su recompensa según el déficit
observado o estimado. La lectura dual linealizada es que $r_k$ funciona localmente como
multiplicador de una restricción de servicio; no hay subasta combinatoria ni asignador
central de coaliciones.

\begin{figure}[H]
\centering
\begin{tikzpicture}[
  block/.style={draw=VIUDark,fill=VIULight,rounded corners=2pt,
    align=center,minimum width=3.0cm,minimum height=0.85cm,font=\small},
  arrow/.style={-{Latex[length=2mm]},thick,color=VIUOrange!90!black}
]
\node[block] (r) {precio\\$r_k$};
\node[block, right=0.8cm of r] (f) {payoff\\$F_{ik}$};
\node[block, right=0.8cm of f] (y) {Smith\\$y_i$};
\node[block, below=0.85cm of y] (z) {capacidad\\$\Zeff_k$};
\node[block, below=0.85cm of r] (e) {error\\$q_k-\Zeff_k$};
\draw[arrow] (r) -- (f);
\draw[arrow] (f) -- (y);
\draw[arrow] (y) -- (z);
\draw[arrow] (z) -- (e);
\draw[arrow] (e) -- (r);
\end{tikzpicture}
\caption{Subasta evolutiva de un reloj. El precio no asigna robots directamente; cambia
el gradiente que ven las dinámicas de Smith.}
\label{fig:marco_subasta}
\end{figure}

\begin{teorema}[Precio de equilibrio] \textbf{[PROBADO+GATE]}
Si la tarea $k$ es factible y el objetivo de ocupación es $q_k$, el precio que deja al
robot marginal indiferente frente a inactividad es
\[
  r_k^\dagger=\frac{r_0}{\Phi_k(q_k)}.
\]
Si la tarea es infactible dentro del intervalo de precios, la proyección satura
$r_k=r_{\max}$ y el reparto resultante se interpreta como water-filling de escasez.
\end{teorema}

\begin{proof}
En equilibrio factible se tiene $\Zeff_k=q_k$. El robot marginal compara tarea con
inactividad, de modo que $r_k^\dagger\Phi_k(q_k)=r_0$. Despejar da la expresión. Si no
existe precio admisible que alcance $q_k$, la ley proyectada alcanza el borde superior
y las ecuaciones de equilibrio vuelven al caso de escasez.
\end{proof}

\begin{proposicion}[Subasta como dual linealizado] \textbf{[PROBADO+GATE]}
Considérese la restricción $G_k(\Zeff_k)\ge G_k(q_k)$ y el Lagrangiano
\[
  \mathcal{L}
  =
  r_0\sum_i y_{i0}
  +
  \sum_k\lambda_k\,[G_k(\Zeff_k)-G_k(q_k)].
\]
Entonces
\[
  \frac{\partial \mathcal{L}}{\partial y_{ik}}
  =
  \lambda_k\Phi_k(\Zeff_k)s_{ik}.
\]
Con $\lambda_k=r_k$, el payoff es el gradiente primal de la restricción de servicio.
\end{proposicion}

\begin{proof}
La prueba es regla de la cadena:
\[
  \frac{\partial \mathcal{L}}{\partial y_{ik}}
  =
  \lambda_kG_k'(\Zeff_k)
  \frac{\partial \Zeff_k}{\partial y_{ik}}
  =
  \lambda_k\Phi_k(\Zeff_k)s_{ik}.
\]
Además, el ascenso dual exacto
$\dot\lambda_k=\gamma[G_k(q_k)-G_k(\Zeff_k)]$ y la ley
\eqref{eq:marco_subasta} comparten punto fijo. Cerca de $\Zeff_k=q_k$,
\[
  G_k(q_k)-G_k(\Zeff_k)
  =
  \Phi_k(q_k)(q_k-\Zeff_k)+O((q_k-\Zeff_k)^2),
\]
por lo que las tasas locales sólo difieren por el factor positivo $\Phi_k(q_k)$.
\end{proof}

Por tanto, \eqref{eq:marco_subasta} no se presenta como primal--dual exacto. La
afirmación defendible es que actúa como control integral equivalente, en primer orden,
al ascenso dual suavizado alrededor de la restricción activa.

\subsection{Energía maestra}

Para cerrar juego y movimiento se introduce una energía:
\begin{equation}
  E(p,y)=
  -r_0\sum_i y_{i0}
  -\sum_{k=1}^{K}r_kG_k(\zeff_k)
  +U_{\mathrm{rep}}(p)+U_{\mathrm{obs}}(p),
  \label{eq:marco_energia}
\end{equation}
donde
\[
  \zeff_k=\sum_i\Psi(T_k(p_i))y_{ik}.
\]
El uso de ocupación efectiva no es cosmético: si se usa la ocupación cruda, las derivadas
cruzadas con respecto a posición y preferencia dejan de coincidir y se pierde la lectura
de potencial.

La identidad de descenso corresponde al flujo continuo ideal: precios congelados, mapa
fijo, $\Psi$ suave, punto holonómico o punto adelantado perfectamente controlado, sin
saturaciones cinemáticas, sin ORCA externo y con contribuciones $s_{ik}$ independientes
de $y$ al derivar. En la implementación real se usa como estructura guía, no como
garantía global del simulador completo.

\begin{teorema}[Descenso Smith--campo] \textbf{[PROBADO+GATE, flujo ideal]}
Con precios congelados y ocupación efectiva coherente con el payoff, el sistema
\[
  \dot y=\mathrm{Smith}(y,F),
  \qquad
  \dot p=-\kappa\nabla_pE
\]
satisface
\[
  \dot E
  =
  -\mathcal{D}_{\mathrm{Smith}}(y,F)
  -
  \kappa\sum_i\norm{\nabla_{p_i}E}^2
  \le0,
\]
donde $\mathcal{D}_{\mathrm{Smith}}\ge0$.
\end{teorema}

\begin{proof}
La derivada total se separa en dos términos:
\[
  \dot E=\nabla_yE^\top\dot y+\nabla_pE^\top\dot p.
\]
Como $-E$ es el potencial que Smith asciende en $y$, el primer término es no positivo y
se escribe como $-\mathcal{D}_{\mathrm{Smith}}$. La ley de campo da
\[
  \nabla_pE^\top\dot p
  =
  -\kappa\sum_i\norm{\nabla_{p_i}E}^2.
\]
La suma prueba el descenso.
\end{proof}

Con uniciclo saturado, paso discreto, transición híbrida y evitación reactiva de
obstáculos, la monotonía exacta de \eqref{eq:marco_energia} se reemplaza por el lema de
paso discreto bajo suavidad local y por validación empírica.

\begin{corolario}[Diagrama de potencia]
Si $\Psi(d)=\exp(-d^2/(2\lambda^2))$, las tareas tienen posiciones fijas, $\zeff_k$ se
congela al comparar y se ignoran obstáculos o se usa una geodésica fija, la tarea
preferida por el robot $i$ satisface
\[
  k^*(i)=
  \argmin_k
  \left[
    \frac{d_{ik}^2}{2\lambda^2}
    -
    \log(r_k\Phi_k(\zeff_k))
  \right].
\]
El equilibrio espacial induce una partición tipo Laguerre/potencia con pesos endógenos;
no se afirma equivalencia con transporte óptimo completo.
\end{corolario}

\subsection{Teoría heterogénea: curva de oferta de la flota}

La extensión más importante del reporte avanzado es reemplazar conteo de robots por
capacidad efectiva. Supóngase que existen $P$ clases de robots con capacidades
\[
  0<c^1<\cdots<c^P,
\]
y $N^p$ robots de clase $p$. Sea $a_k^p$ la masa de clase $p$ asignada a la tarea $k$.
En esta subsección $a_k^p\in\mathbb{R}_+$ es una relajación fluida; si se impone
integralidad, el problema de clearing entero se convierte en una variante de
\emph{knapsack}.
La capacidad agregada es
\[
  Z_k=\sum_{p=1}^{P}c^p a_k^p.
\]
El potencial agregado es
\begin{equation}
  V(a)=r_0\sum_p a_0^p+\sum_{k=1}^{K}r_kG_k(Z_k).
  \label{eq:marco_potencial_hetero}
\end{equation}

\begin{lema}[Gradiente heterogéneo]
El potencial \eqref{eq:marco_potencial_hetero} satisface
\[
  \frac{\partial V}{\partial a_k^p}=c^p r_k\Phi_k(Z_k).
\]
\end{lema}

\begin{proof}
Por regla de la cadena,
\[
  \frac{\partial V}{\partial a_k^p}
  =
  r_kG_k'(Z_k)\frac{\partial Z_k}{\partial a_k^p}
  =
  r_k\Phi_k(Z_k)c^p.
\]
\end{proof}

\begin{teorema}[Orden de despliegue por capacidad] \textbf{[PROBADO+GATE]}
Si la inactividad paga $r_0$ por robot, en un óptimo no puede haber un robot ligero
activo mientras un robot más pesado permanece inactivo, salvo indiferencia exacta.
\end{teorema}

\begin{proof}
Considérese un robot ligero activo de capacidad $c^p$ y un robot pesado inactivo de
capacidad $c^q>c^p$. Intercambiar sus roles mantiene el número de robots activos, por lo
que el pago de inactividad cambia de forma simétrica, pero incrementa la capacidad en la
tarea en $c^q-c^p$. Como $G_k$ es creciente y cóncava, ese incremento no reduce el
potencial y lo mejora si la tarea conserva valor marginal. Por tanto una asignación con
esa inversión no puede ser óptima.
\end{proof}

La flota induce una función de coste mínimo de desplegar capacidad $Q$:
\[
  h(Q)=
  \min\left\{
    \sum_p u_p:
    0\le u_p\le N^p,\;
    \sum_pc^pu_p\ge Q
  \right\}.
\]
Esta función es convexa y lineal por tramos. Sus pendientes son
\[
  \frac{1}{c^P}<\frac{1}{c^{P-1}}<\cdots<\frac{1}{c^1},
\]
porque los robots pesados son más baratos por unidad de capacidad desplegada cuando la
inactividad paga por robot.
Con $\rho_p=r_0/c^p$, el orden de capacidades implica
\[
  \rho_P<\rho_{P-1}<\cdots<\rho_1.
\]

\begin{figure}[H]
\centering
\begin{tikzpicture}[scale=0.95]
  \draw[-{Latex[length=2mm]},thick] (0,0) -- (8.2,0) node[right] {$Q$};
  \draw[-{Latex[length=2mm]},thick] (0,0) -- (0,5.0) node[above] {$\lambda$};
  \draw[VIUOrange,very thick] (0.4,0.85) -- (2.1,0.85) -- (2.1,1.95) --
    (4.3,1.95) -- (4.3,3.55) -- (7.4,3.55);
  \draw[VIUDark,very thick,domain=0.65:7.3,samples=90]
    plot (\x,{4.35/(1+exp(0.78*(\x-3.75)))+0.42});
  \draw[dashed,VIUGray] (3.15,0) -- (3.15,1.95) -- (0,1.95);
  \fill[VIUDark] (3.15,1.95) circle (2pt);
  \node[anchor=north] at (3.15,-0.05) {$Q^*$};
  \node[anchor=east] at (-0.05,1.95) {$\lambda^*$};
  \node[VIUOrange,anchor=west,font=\small] at (4.65,3.85) {oferta de flota};
  \node[VIUDark,anchor=west,font=\small] at (4.7,1.05) {demanda de tareas};
  \draw[decorate,decoration={brace,amplitude=5pt},VIUGray]
    (2.1,-0.28) -- (4.3,-0.28)
    node[midway,below=6pt,font=\small] {kink posible};
\end{tikzpicture}
\caption{Mercado de capacidad efectiva. La oferta es una escalera inducida por la flota;
la demanda proviene de invertir las sigmoidales de servicio.}
\label{fig:marco_mercado}
\end{figure}

La demanda agregada de capacidad al precio sombra $\lambda$ se escribe de nuevo por
tramos para respetar el dominio del logaritmo:
\[
  W_k^M(\lambda)=
  \begin{cases}
  0,
  & \lambda\ge r_k\Phi_k(0),\\[4pt]
  M_k+\dfrac{1}{\beta}\log\!\left(\dfrac{r_k}{\lambda}-1\right),
  & 0<\lambda<r_k\Phi_k(0),
  \end{cases}
\]
y
\[
  D(\lambda)=
  \sum_k W_k^M(\lambda).
\]

\begin{teorema}[Water-filling con curva de oferta] \textbf{[PROBADO+GATE, modelo fluido agregado]}
El equilibrio heterogéneo agregado queda caracterizado por
\[
  \lambda^*\in\partial(r_0h)(Q^*),
  \qquad
  Q^*=D(\lambda^*),
  \qquad
  Z_k^*=W_k^M(\lambda^*).
\]
El precio puede estar en un tramo de oferta, en un kink de población agotada o en
escasez total.
\end{teorema}

\begin{proof}
Por el teorema de orden de despliegue, el problema agregado puede reducirse a elegir
cuánta capacidad total desplegar, pagando coste $r_0h(Q)$. La condición de optimalidad
es $\lambda^*\in\partial(r_0h)(Q^*)$. Para cada tarea activa, la condición marginal
$r_k\Phi_k(Z_k^*)=\lambda^*$ da la fórmula de $Z_k^*$ al invertir la sigmoide. La suma
de demandas por tarea cierra $Q^*=D(\lambda^*)$.
\end{proof}

El resultado anterior es un clearing fluido de capacidad. El paso a robots enteros
requiere redondeo, histéresis o una rutina discreta de selección; esa capa no está
incluida en la prueba del mercado continuo.

\begin{observacion}[Kink multi-tarea]
Cuando una clase de robots se agota, $\lambda^*$ no tiene por qué coincidir con
$r_0/c^p$ para ninguna clase. En un problema con varias tareas, el precio en el kink
resuelve $\sum_k W_k(\lambda)=A_p$. La fórmula de una sola tarea evaluada sobre la
capacidad total sólo es válida para $K=1$.
\end{observacion}

\subsection{Presupuesto de percepción}

La robustez del sistema depende de una cadena de errores: sensor, localización, ruta,
consenso, payoff, ocupación y redondeo entero. El resultado útil para la memoria es que
el redondeo entero protege si el error permanece dentro de una zona muerta.

\begin{figure}[H]
\centering
\begin{tikzpicture}[
  block/.style={draw=VIUDark,fill=VIULight,rounded corners=2pt,
    align=center,minimum width=2.35cm,minimum height=0.78cm,font=\small},
  arrow/.style={-{Latex[length=2mm]},thick,color=VIUOrange!90!black}
]
\node[block] (a) {sensor\\$\sigma_m$};
\node[block, right=0.35cm of a] (b) {pose\\$P_\infty$};
\node[block, right=0.35cm of b] (c) {ruta\\$T,\Psi$};
\node[block, right=0.35cm of c] (d) {DAC\\$e_Z$};
\node[block, below=0.82cm of d] (e) {payoff\\$\bar e$};
\node[block, left=0.35cm of e] (f) {sensibilidad\\$S_k$};
\node[block, left=0.35cm of f] (g) {margen\\$s/2$};
\node[block, left=0.35cm of g] (h) {coalición\\misma};
\draw[arrow] (a)--(b);
\draw[arrow] (b)--(c);
\draw[arrow] (c)--(d);
\draw[arrow] (d)--(e);
\draw[arrow] (e)--(f);
\draw[arrow] (f)--(g);
\draw[arrow] (g)--(h);
\end{tikzpicture}
\caption{Cadena del presupuesto de percepción: un error físico sólo cambia la coalición
si logra cruzar el margen entero de redondeo.}
\label{fig:marco_percepcion}
\end{figure}

\begin{teorema}[Presupuesto de percepción] \textbf{[PROBADO+GATE]}
Sea $\bar e$ una cota uniforme de error de payoff y
\[
  S_k=\frac{(1+e^\beta)^2}{\beta r_ke^\beta}
\]
la sensibilidad local de ocupación en $q=n+1$. Si
\[
  \max_k S_k\bar e < \frac{s}{2},
\]
donde $s$ es el margen entero de redondeo, entonces el clearing selecciona las mismas
coaliciones que el sistema nominal.
\end{teorema}

\begin{proof}
La inversa local del payoff marginal da
\[
  \delta z_k=
  \frac{\delta f_k}{\beta r_k\Phi_k(1-\Phi_k)}.
\]
En $q=n+1$, sustituir $\Phi(q)=1/(1+e^\beta)$ produce $S_k$. Si el desplazamiento de
ocupación es menor que la mitad del margen entero, ninguna tarea cruza el umbral de
redondeo que define la coalición.
\end{proof}

La prueba de Lyapunov continua de Smith no se transfiere automáticamente a la versión
Euler-proyectada. En el simulador, la proyección al simplex es una salvaguarda numérica;
la estabilidad discreta queda condicionada a paso suficientemente pequeño y a la
suavidad local de la energía.

\subsection{Tiempo discreto e híbrido}

El simulador no ejecuta un flujo continuo, sino un bucle discreto. Por tanto, el
descenso de energía necesita una condición de paso.

\begin{lema}[Paso de Euler por suavidad] \textbf{[CONDICIONADO]}
Si $E$ es $L$-suave en la región visitada y
\[
  x^+=x-\eta\nabla E(x),
\]
entonces
\[
  E(x^+)\le E(x)-\eta\left(1-\frac{L\eta}{2}\right)\norm{\nabla E(x)}^2.
\]
Hay descenso si $0<\eta<2/L$.
\end{lema}

\begin{proof}
Por suavidad,
\[
  E(x^+)\le E(x)+\nabla E(x)^\top(x^+-x)+\frac{L}{2}\norm{x^+-x}^2.
\]
Sustituyendo $x^+-x=-\eta\nabla E(x)$ se obtiene la desigualdad.
\end{proof}

\begin{proposicion}[Conmutación finita del clearing entero] \textbf{[CONDICIONADO]}
Si una transición entera de coalición sólo se acepta cuando $\Delta V>\varepsilon$ y
$V\in[V_{\min},V_{\max}]$, entonces
\[
  N_{\mathrm{switch}}\le\frac{V_{\max}-V_{\min}}{\varepsilon}.
\]
\end{proposicion}

\begin{proof}
Cada transición aceptada consume al menos $\varepsilon$ de incremento de potencial. La
suma total de incrementos no puede exceder $V_{\max}-V_{\min}$.
\end{proof}

\subsection{Extensiones avanzadas}

Las siguientes ideas son útiles como cierre del TFM y como puente doctoral, pero no deben
presentarse con el mismo estado que los teoremas anteriores.

\subsubsection{Mercado de supervivencia por batería}

La recarga puede formularse como una tarea adicional:
\[
  F_{i,\mathrm{ch}}
  =
  R_{\mathrm{ch}}\,
  \sigma(\alpha(B_{\mathrm{safe}}-B_i))\,
  \Psi(T_{\mathrm{ch}}(p_i)).
\]
La afirmación correcta es una condición de dominancia, no una barrera infinita: una
sigmoide está acotada.

\begin{lema}[Dominancia de recarga] \textbf{[FUTURO]}
Si para $B_i<B_{\mathrm{umbral}}$ se cumple
\[
  F_{i,\mathrm{ch}}>\max_k F_{ik}+\delta,
  \qquad \delta>0,
\]
entonces Smith transfiere masa hacia recarga con tasa inferior proporcional a $\delta$
mientras haya masa en tareas ordinarias.
\end{lema}

\begin{proof}
Para cada tarea ordinaria $k$ con $y_{ik}>0$,
\[
  \pos{F_{i,\mathrm{ch}}-F_{ik}}\ge\delta.
\]
El término de entrada hacia recarga en Smith contiene la suma de esos gaps ponderados
por $y_{ik}$.
\end{proof}

\subsubsection{Payoffs predictivos}

Un payoff predictivo puede escribirse como
\[
  J_{ik}^{H}
  =
  \max_u
  \int_t^{t+H}
  \left[
    r_k\Phi_k(\hat s_{ik}(\tau))-\ell_B(B_i(\tau),u(\tau))
  \right]\dd\tau
  +V_T.
\]
El problema es que, en general, $J^H$ rompe la simetría de derivadas cruzadas del juego
potencial. La ruta honesta es tratarlo como perturbación:
\[
  |J_{ik}^{H}-F_{ik}|\le\bar\delta.
\]
Si el núcleo nominal tiene una cota ISS con radio $C\bar e$, el predictivo tendría radio
$C(\bar e+\bar\delta)$.

\subsubsection{Amortiguamiento anticipatorio en coordenadas agregadas}

La anticipación de primer orden usa
\[
  F^\tau=F+\tau\dot F.
\]
Si $F=\nabla V$, entonces
\[
  F^\tau=\nabla V+\tau\nabla^2V\,\dot y.
\]
El punto delicado es que $\nabla^2V$ no es definida negativa de rango completo. Para
\[
  V(y)=\sum_k r_kG_k\left(\sum_i y_{ik}\right),
\]
el Hessiano tiene rango a lo sumo $K$, porque $V$ sólo depende de las ocupaciones
agregadas. El amortiguamiento anticipatorio actúa en el subespacio $Z$, no en todas las
redistribuciones microscópicas de $y$. El término técnico \emph{Hessian damping} se usa
sólo en ese sentido restringido.

\begin{lema}[Rango agregado del Hessiano] \textbf{[FUTURO con primer gate]}
El Hessiano de $V(y)=\sum_k r_kG_k(\sum_i y_{ik})$ respecto a $y$ tiene rango a lo sumo
$K$. Las direcciones que preservan todas las ocupaciones $Z_k$ pertenecen al núcleo.
\end{lema}

\begin{proof}
Sea $A$ el operador lineal que lleva $y$ a $Z=Ay$. Entonces
$V(y)=\bar V(Ay)$ y
\[
  \nabla_y^2V=A^\top\nabla_Z^2\bar V\,A.
\]
Por tanto $\operatorname{rank}(\nabla_y^2V)\le K$. Si $A\delta y=0$, entonces
$\delta y$ no cambia ninguna ocupación agregada y queda en una dirección nula.
\end{proof}

\subsubsection{Eikonal predictiva y protocolo de jets}

La eikonal predictiva reemplaza la distancia actual por una velocidad admisible
deformada por densidad futura:
\[
  \rho(x,t)=\sum_i K_h(x-p_i(t)),
  \qquad
  \rho^\tau=\clip\!\left[\rho+\tau\dot\rho+\frac{\tau^2}{2}\ddot\rho\right],
\]
\[
  v^\tau(x)=v_{\min}+(v_{\max}-v_{\min})e^{-\alpha\rho^\tau(x)},
  \qquad
  \norm{\nabla T_k^\tau(x)}=\frac{1}{v^\tau(x)}.
\]
La regla de comunicación compatible con el enfoque distribuido es no compartir rutas
completas, sino un jet local:
\[
  \mathcal{J}_i=(p_i,u_i,y_i,\dot y_i,B_i,\dot B_i,\nabla\mathcal{U}_i,
  \nabla^2\mathcal{U}_i).
\]
Con él, un vecino aproxima
\[
  p_i(t+\tau)\approx p_i(t)+\tau u_i(t)+\frac{\tau^2}{2}a_i(t).
\]
Esto conserva el principio de arquitectura: se comparte el mundo, no la voluntad; la
intención sólo aparece como derivadas locales de corto horizonte.
Como $T_k^\tau$ cambia con la densidad prevista y, por tanto, con el tiempo, esta
extensión rompe la energía estática \eqref{eq:marco_energia}. Debe mantenerse como
\textbf{[FUTURO]} y no mezclarse con la prueba de $\dot E\le0$.

\subsection{Síntesis de resultados y uso en la memoria}

\begingroup
\small
\setlength{\tabcolsep}{3pt}
\begin{longtable}{@{}p{0.20\linewidth}p{0.15\linewidth}p{0.17\linewidth}p{0.18\linewidth}p{0.13\linewidth}@{}}
\toprule
\textbf{Resultado} & \textbf{Ecuación clave} & \textbf{Estado} & \textbf{Gate o evidencia} & \textbf{Uso} \\
\midrule
Water-filling homogéneo & \eqref{eq:marco_water_hom} & PROBADO+GATE &
\texttt{verify\_teoria} & Cuerpo \\
Perturbación DAC & Prop.~\ref{prop:perturbacion} & PROBADO local &
Escenario D pendiente & Cuerpo \\
Energía maestra ideal & \eqref{eq:marco_energia} & PROBADO+GATE ideal &
\texttt{verify\_energia} & Breve \\
Mercado heterogéneo fluido & \eqref{eq:marco_potencial_hetero} & PROBADO+GATE fluido &
gate de extensiones & Anexo/cuerpo \\
ISS práctica & \eqref{eq:anexo_iss_limsup} & CONDICIONADO &
pendiente & Anexo \\
ORCA económico & \eqref{eq:anexo_alpha_orca} & CONDICIONADO &
pendiente & Futuro \\
Retornabilidad y relevo & \eqref{eq:anexo_retornabilidad} & FUTURO &
pendiente & Futuro \\
Anticipatorio muestreado & \eqref{eq:anexo_descenso_sampled} & FUTURO &
pendiente & Futuro \\
Densidad/eikonal predictiva & \eqref{eq:anexo_densidad_ddot} & FUTURO &
pendiente & Futuro \\
Benchmark T1--T5 & Sec.~\ref{sec:resultados} & PENDIENTE DE CIERRE &
v2.7 y CoppeliaSim & Cuerpo final \\
\bottomrule
\end{longtable}
\endgroup

\begingroup
\small
\setlength{\tabcolsep}{3pt}
\begin{longtable}{@{}p{0.24\linewidth}p{0.66\linewidth}@{}}
\toprule
\textbf{Resultado} & \textbf{Hipótesis esenciales que deben acompañarlo} \\
\midrule
Teorema homogéneo & $\Psi\equiv1$, comunicación global, robots homogéneos,
sin transición de modo, balance exacto $N=Kn$ y paso pequeño si se discretiza. \\
Energía maestra & $s_{ik}$ no depende de $y$ al derivar, mapa fijo, flujo continuo,
sin saturación, sin ORCA externo y sin mínimos locales usados como garantía global. \\
Mercado heterogéneo & Relajación fluida $a_k^p\in\mathbb{R}_+$, inactividad pagada por
robot, $G_k$ creciente, capacidades ordenadas y clearing entero tratado aparte. \\
ISS práctica & Decaimiento nominal local, error de payoff acotado y equivalencia local
entre $W$ y la distancia al conjunto de Nash. \\
ORCA económico & $w_i>0$, factibilidad par-a-par, semiplanos ORCA con intersección
no vacía o relajación penalizada/parada segura. \\
Anticipatorio & Horizonte $\tau$ pequeño, suavidad local, error de muestreo acotado y
amortiguamiento sólo en coordenadas agregadas. \\
Eikonal predictiva & Densidad suavizada, claridad geodésica local y separación explícita
respecto a la prueba de energía estática. \\
\bottomrule
\end{longtable}
\endgroup

El Anexo~\ref{sec:anexo-teorico-extendido} recoge las ecuaciones extendidas que no
conviene insertar completas en el cuerpo principal, pero que sí son útiles para auditoría
formal: lift general de contribución efectiva, condición inversa de integrabilidad,
curva de oferta por tramos, ISS explícita, ORCA económico, retornabilidad energética,
handover dinámico, Smith anticipatorio muestreado y densidad predictiva de segundo
orden. La recomendación de escritura para la memoria final es usar este capítulo como
soporte teórico y no como reemplazo del benchmark. La contribución debe defenderse en
tres capas: primero, el núcleo homogéneo probado; segundo, la validación experimental
del bucle completo; tercero, la extensión heterogénea y predictiva como frontera
avanzada con estado declarado.
